% A dictionary between finite spectral triples and Kaluza-Klein data: short note, synthesis of lemmas 05 to 30 and the 11D line
\ifdefined\mobile
\documentclass[aps,prd,reprint,onecolumn,nofootinbib,amsmath,amssymb,floatfix]{revtex4-2}
\usepackage[paperwidth=270pt,paperheight=585pt,textwidth=246pt,textheight=545pt,top=18pt,left=12pt,headheight=0pt,headsep=0pt,footskip=14pt]{geometry}
\else
\documentclass[aps,prd,reprint,nofootinbib,amsmath,amssymb,floatfix]{revtex4-2}
\fi
\usepackage[T1]{fontenc}
\usepackage{lmodern}
\usepackage{amsthm}
\emergencystretch=1.5em
\ifdefined\mobile\makeatletter\def\frontmatter@abstractwidth{246\p@}\makeatother\fi
\usepackage[hidelinks]{hyperref}
\theoremstyle{plain}
\newtheorem{proposition}{Proposition}
\newcommand{\C}{\mathbb C}
\renewcommand{\H}{\mathbb H}
\newcommand{\CP}{\mathbb{CP}}
\newcommand{\supp}[1]{Supplemental Material, part #1}
\newcommand{\file}[1]{Supplemental Material, 11D file #1}

\begin{document}

\title{A dictionary between finite spectral triples and Kaluza-Klein data}
\date{Draft, 4 November 2026}

\begin{abstract}
We collect, with the status of each, the statements that translate Connes' finite spectral triple into Kaluza-Klein data. The grading of the two-point space is the reflection of an antiperiodic circle. The finite Dirac operator is a wall angle and a flat connection. The finite Higgs one-form is the compression of the gauge field along the internal direction, and the Higgs modulus is the angle between the wall gradings. The real structure decides whether a massless block has orthogonal or symplectic symmetry, which fixes the antisymmetric or symmetric diquark block. Noncommutative algebras arise as windows of commutative parents. The translation has limits, which we state: the finite data and the chirality remain inputs on the interval; the window algebras concern compression maps; no action in the flat family selects a path on which a dimension opens, and with the top alone it has no vacuum in the window. A reading through the heterotic and type I$'$ web fails in three specific ways: the angle-position map is nonlinear, the length of the interval is T-dual to the type I$'$ length, and the relevant configuration has a Romans mass, with no eleven-dimensional regime.
\end{abstract}

\maketitle

\section{Introduction}

Connes' finite spectral triple packages the masses, the Higgs field and the gauge group of the Standard Model into finite-dimensional data. A Kaluza-Klein theory packages them into the geometry of internal dimensions. The proposal of \cite{towerlesspaper} is that the first is a towerless window of the second, with 9 and 11 as limits of an even interior. This note asks what the word ``is'' can mean here, entry by entry. We separate what is proved from what is a reading, and we say where the translation stops. It brings together the two companion notes \cite{boundarynote,jointnote} and a line of work on eleven-dimensional completions. It contains few new statements of its own; those are marked.

Table~\ref{tab:dictionary} is the summary. The status column uses the record in the repository ledger: ``refereed'' means an independent reading by GPT-6.1 Terra or GPT-6 Astra, ``one reading'' an independent reading by Claude Sonnet 5.5, and ``unrefereed'' a statement checked only by its author.

\begin{table*}
\caption{\label{tab:dictionary}The dictionary. Parts refer to the Supplemental Material; 11D/n is file n of the eleven-dimensional line.}
\begin{ruledtabular}
\begin{tabular}{p{0.27\textwidth}p{0.47\textwidth}p{0.18\textwidth}}
Finite side&Kaluza-Klein side&Status\\
\hline
Two points, grading $\gamma$&Antiperiodic circle of radius $l/\pi$; $\gamma$ is its reflection, the walls its fixed points&Proposition, 11D/01, lemma 05\\
$D_F$ with KO-6 signs&Wall data $X=\tan(lD_F)$, with exact window; a constant connection, odd for the grading&Theorem, refereed, lemma 22\\
$J_F$, $\gamma_F$&Exact on the parent&Theorem, refereed, lemma 22\\
Masses&Internal momenta: $(n\pi\pm\theta_k)/l$, $\theta_k=l\sigma_k$&Theorem, refereed, lemma 22\\
Higgs one-form $[D_F,a]$&Compression of the connection along the interval; the modulus is the angle between two wall gradings&Theorems, refereed, lemmas 22, 23; the interpretation has one reading, 11D/11\\
Gauge fields, centralizer $C_F$&$A_\mu$ with wall conditions; a generator has a massless vector or a massless scalar, never both&Proposition 22.8, refereed; Theorem 11.3, one reading\\
Sign $s$ of an odd factorization&$j^2=s$: $O(n)$ or $USp(n)$ on a massless block; $SO(30)$ of the five light quarks for $s=+1$&Propositions, 11D/08 unrefereed, 11D/11 one reading\\
Coefficients $(3,24/5)$ of the action&$(2,16/5)$ of the interval, joined at leading order by $(3-d_{\rm eff})/2$&Theorem, refereed, lemma 25\\
$M_2(\H)\oplus M_4(\C)$&Limit of compressions of real functions on a disc bundle over $S^4$&Theorems, refereed, lemmas 28, 29\\
$M_p(\C)$, KO-6&Lowest Landau level of a doubled torus with flux $p$, with the unital datum $M_p(\C)\oplus\C$&Proposition, refereed with amendment, 11D/07\\
$M_W,M_Z\to0$&A family with $m\ll E_\varepsilon\ll\Delta\to0$ and a divergent count at fixed energy: vector poles in a five-dimensional toy, and a fermion window only on the disc bundle&Theorems, refereed, lemma 30\\
\end{tabular}
\end{ruledtabular}
\end{table*}

\section{The grading is a reflection}

Let $D_l(L,R)=(R',-L')$ on $L^2([0,l],\C^2)$, with $R(0)=0$ and $L(l)=0$, grading $\Gamma=\mathrm{diag}(1,-1)$ and scalar functions on $[0,l]$: the chiral interval whose minimal truncation is the Connes-Lott two-point space \cite{ConnesLott1991}. Then
\begin{equation}
(W(L,R))(x)=\tfrac1{\sqrt2}\bigl(L(|x|)+i\,\mathrm{sgn}(x)R(|x|)\bigr)
\end{equation}
is a unitary onto the antiperiodic functions on a circle of length $2l$, with $WD_lW^*=-i\,d/dx$, $W\Gamma W^*$ the reflection $x\to-x$, and the algebra going to the reflection-even functions (\file{01}). So the two sheets of the finite geometry are the two fixed points of a reflection, and the grading of the finite geometry is that reflection. The Connes-Lott mass is $m=\pi/(2l)=1/(2R)$, the lowest momentum of an antiperiodic spinor. The two-point distance $1/m$ is a statement about this circle.

\section{The finite Dirac operator is a boundary condition and a holonomy}

Every $D_F$ with the signs of KO-dimension six is exactly the lowest window of the pure derivative on the interval, for $l\|D_F\|<\pi/2$, with the wall data $X=\tan(lD_F)$, and with $J_F$ and $\gamma_F$ exact. There is no bulk mass, and the masses are momenta. After a unitary change of domain, $D_F$ is a constant connection, odd for the grading \cite{boundarynote}. On the periodic circle of \supp{25} it is a Wilson line, and the angle $\theta=l\sigma$ of a channel lives on the fundamental domain $[0,\pi/2]$, whose ends are the periodic and the antiperiodic spectra (\file{07}).

The holonomy form has a consequence for dynamics. A flat connection cannot be seen by a local operator in the bulk. So the dependence of an action on $D_F$ at large $l$ is nonlocal, of the Casimir kind, and it falls with $l$; this is the point of Section~\ref{sec:path}.

\section{The Higgs field is the gauge field along the interval}

In Connes' geometry the Higgs field comes from inner fluctuations $\sum a_i[D_F,b_i]$. In a Kaluza-Klein theory a scalar in the gauge sector comes from the component of a connection along the internal direction. That the Higgs is a connection in the discrete internal direction is the original reading of \cite{ConnesLott1991}. With a parent, it becomes a statement relating an operator on the interval to an operator on the window, in three precise senses.

First, for $f$ in the parent algebra, the compression $Q_l$ to the window satisfies $Q_l([D_X,f])=[D_F,Q_l(f)]$ (\supp{22}). More generally $Q_l(ET)=-l^{-1}\int_0^l\sin(x\,\mathrm{ad}_{D_F})(T(x))\,dx$, which is $-\tfrac l2[D_F,T]$ to leading order for constant $T$. So a connection along the interval compresses to a finite one-form.

Second, a connection $h$ along the interval is equivalent to a rotation of the second wall: $D(\Gamma_0,\Gamma_1)+Eh=uD(\Gamma_0,w^*\Gamma_1w)u^*$, with $w$ its holonomy (\supp{23}). The fluctuations odd for the grading change the principal angles between the two wall gradings, which are the masses. So the Higgs modulus is the angle between the walls. For $M_2(\H)$ the field space is $S^4$, and the modulus is half the geodesic distance between the points chosen by the walls. The even fluctuations, by contrast, leave all the singular values fixed (\supp{22}, No-go theorem 22.7).

Third, which of these are physical is decided at the walls. With wall algebras of skew parts $k_0,k_l$ and a five-dimensional Yang-Mills action, the massless scalars from the interval component are the orthogonal complement of $k_0+k_l$. A generator in $k_0\cap k_l$ has a massless vector and no massless scalar; a generator in the complement has a massless scalar and no massless vector; all others have neither (\file{11}). So in the parent of lemma 22, whose algebra is free at one wall, the interval component has no zero mode, and the Higgs is the wall datum $X$. In the two-grading parent of lemma 23, it is a modulus. The rule ``vector or scalar, never both'' is the reason a gauge-Higgs scalar in a chosen representation forces the walls.

\section{The real structure fixes the symmetry of a massless block}

An odd factorization of a KO-6 triple uses a grading-reversing unitary $U$ with $UJ=sJU$; $s=+1$ gives KO-7 (label 11) and $s=-1$ gives KO-5 (label 9) \cite{towerlesspaper}. The induced antiunitary $j=JU$ on the positive-chirality subspace has $j^2=s$, so the unitaries commuting with it on a massless block of dimension $n$ form $O(n)$ for $s=+1$ and $USp(n)$ for $s=-1$ (\file{08}).

\begin{proposition}\label{prop:so30}
For the five light quarks in three colours, the positive-chirality block is $H_+=\{q_L\}\oplus\{\overline{q_R}\}=V\oplus\bar V$, of dimension 30, with $j\,q_L=s\,\overline{q_R}$ and $j\,\overline{q_R}=q_L$. For $s=+1$ the skew generators commuting with $j$ are $\mathfrak{so}(30)=\mathfrak u(15)\oplus\Lambda^2V$, with $\mathfrak u(15)$ the commutant of quark number, and for $s=-1$ they are $\mathfrak{usp}(30)=\mathfrak u(15)\oplus S^2V$. In the colour antitriplet, with charges $(-2/3,+1/3,+4/3)$, the counts are $(6,6,3)$ and $(3,6,1)$.
\end{proposition}

The proof is linear algebra (\file{11}). So the $SO(30)$ met by counting in the spectral action \cite{towerlesspaper} and the scalar counts of \cite{Rivero2005} have a precise place in the finite geometry: the symmetry of the massless light block under the real structure, for the sign that the seesaw sector also selects. Here $U$ commutes with the vectorlike algebra of label 9, and it serves only to define $j$ for the larger algebra below.

A second statement separates the two kinds of operator. A Hermitian operator $X$ with $JX=\epsilon'XJ$, in the labels 9 to 11, where $\epsilon\epsilon'=+1$, has a symmetric block between particles and antiparticles of one multiplet, as a Majorana mass does (\file{09}). A Hermitian gauge generator anticommutes with $j$ and has an antisymmetric block for $s=+1$ (symmetric for $s=-1$). The two are joined by the commutator with the masses: the connection along the interval has a block $B_{ij}$ in $\Lambda^2V$, and its compression to the window, to first order and with the compression fixed, is the symmetric block
\begin{equation}\label{eq:shadow}
\delta\Lambda_{ij}=-\frac{1-\cos(l(m_i-m_j))}{l(m_i-m_j)}\,B_{ij},
\end{equation}
with $\Lambda$ the mass matrix. This maps $(S^25,\bar3)$ onto $(\Lambda^25,\bar3)$ for distinct masses, with kernel the five same-flavour diquarks, and it vanishes for a universal mass, where $M=mU$ is invariant under $SO(30)$ (\file{11}). So the counts $(6,6,3)$ belong to the parent, and $(3,6,1)$ to its window.

The scope matters. With this algebra order one fails on the block for nonzero masses, and on the interval massless scalars in $\Lambda^2V$ need both walls reduced to $M_{15}(\C)$, as an input. Then the $(24,1)+(1,1)$ are vectors, so the slepton counts of \cite{Rivero2005} do not come out as scalars on an interval.

\section{The algebra}

A function algebra is commutative; a window of it need not be. Three exact examples, in increasing strength. The two lowest levels of the chiral interval give a four-dimensional operator system generating $M_2(\C)$ on each chiral block \cite{towerlesspaper}. The lowest Landau level of a torus with flux $p$, doubled, has the signs of KO-dimension six and, with the unital datum $M_p(\C)\oplus\C$ of \supp{16}, Corollary 16.8, window $M_p(\C)$, with the Landau gap $\sim(p/A)^{1/2}$ sending the tower to infinity as the area shrinks at fixed flux (\file{07}); at $p=3$ this has the form of the colour sector, with the opposite algebra acting on the antiquarks. On the disc bundle of $\Lambda^2_-T^*S^4$, with a Callias potential and a local wall, there are exactly eight light states, a copy of the $\mathbf 4$ of $Sp(2)$ near each end, and real functions compress to maps generating $M_2(\H)\oplus M_4(\C)$ \cite{jointnote}.

Two cautions come with these. Generated algebras need not converge when compression maps do: an unrestricted generated algebra can retain extra directions, and the quaternionic statement holds for bounded expressions, up to a defect (\supp{30}). Second, order zero with the inherited $J$ forces a commutative algebra under pointwise reality (\supp{16}, Proposition 16.1), so a noncommutative window needs a doubling, a second sheet, or a loss of the inherited symmetry; on the disc bundle, the wall breaks charge conjugation, and this is what allows $M_4(\C)$ there.

\section{The massless limit and what selects the path}\label{sec:path}

A path to $M_W=M_Z=0$ opens a dimension only if the first discarded level also falls, more slowly than the light masses, and the count at fixed energy diverges \cite{jointnote}. Such families exist on the interval and on the disc bundle, where the count has seven-dimensional growth. No action has selected one.

A one-loop potential with $l$ and the wall angles as its only scales has the form $l^{-4}G(\theta)$. With $\theta$ free the length runs away, and with $\theta=lm$ an isolated stationary angle fixes $lm$; a flat direction would select nothing (\file{01}). In the normalization of \supp{24}, the angle-independent part is new here:
\begin{equation}\label{eq:a0}
8\pi^2l^4V(0)=\frac{3\zeta(5)}{16}\Bigl(4n_f-3\bigl(n_i-\tfrac{15}{16}n_h\bigr)-n_s-5n_g\Bigr),
\end{equation}
with $n_f$ twice the Dirac species, $n_i$ and $n_h$ the gauge generators with integer and half-integer towers at $\theta=0$, $n_s$ the real scalars with integer towers and $n_g=1$ with a five-dimensional graviton, assuming the bulk vacuum energy and the wall tensions tuned away (\file{10}). For three generations and $\mathfrak{sp}(2)\oplus\mathfrak{su}(4)$, with $n_i=16$ and $n_h=9$, the bracket is $192-22.7=169.3$, so at $v=0$ the length runs to infinity, in the direction of the proposal.

With only the top at the wall angle, $G>0$ on $[0,\pi/2]$, and the length grows until the top leaves the window (Python). With a universal mass, $G$ changes sign and there is a stable point, at $\theta_*=1.009$ for three generations and $\theta_*=1.031$ for one (Python). These angles are of order one, so these vacua do not separate the window. A scale-free action can separate it, with a fixed small ratio, if the light mass is exponentially small in a stabilized shape parameter, as on the disc bundle. A ratio that tends to zero needs a second scale, and a term in $v$ falling more slowly than $m^2/l^2$. Since $D_F$ is a holonomy, the flat family has exactly $m^2/l^2$ at large $l$, so that term must come from elsewhere, naturally a local bulk energy depending on the electroweak order parameter, which goes beyond $D_F$ as a pure boundary condition.

\section{What the string dualities give}

The shape of the heterotic web, with a non-chiral theory in nine dimensions, chiral theories in ten and walls in eleven, suggests a dictionary: the two sheets as two orientifold walls, $D_F$ as sixteen D8-brane positions, $SO(30)\times SO(2)$ as fifteen branes and one displaced for the top \cite{towerlesspaper,PW1996}. At the level of free spectra the top channel is an $SO(32)$ Wilson-line channel, $m_t=a/R$, which is a form of the Hosotani mechanism \cite{Hosotani1983}. Against the sources, the rest needs three corrections (\file{10}).

First, the angle is nonlinear in the position: $\theta_k=\tfrac\pi2W(x_k)/W(L)$ with $W=\int z^{-1/3}dx$, from the BPS masses \cite{PW1996}, so the dilaton is a trace of a function of $|D_F|$ only implicitly. Second, the tower is a winding tower of type I$'$, so the length $l$ and the type I$'$ length $L_p$ are T-dual, with $lL_p=\pi^2\alpha'$ at constant dilaton in the type I$'$ metric, and an extra factor $\lambda_h^{1/2}$ for the heterotic radius in the heterotic metric. The limit $l\to\infty$ at fixed heterotic coupling shrinks the type I$'$ interval, and it ends at ten-dimensional $SO(32)$. Third, the configuration $15+1$ has a Romans mass in the bulk. Its proper length is linear in the inverse coupling, so in the classical profile no region has both strong coupling and a length resolved in eleven-dimensional units, in line with the remark of \cite{BGL1998} that only eight D8-branes with an orientifold plane are known to arise in M-theory.

The charged states massless at the far wall at infinite coupling are D0-branes \cite{BGL1998}, heterotic winding states: they are massless at the string-scale radius and grow heavier with $R$, the opposite of $l\sim1/M_W$. A neutral $Z$ cannot be of this kind, since the rank is fixed. The dilaton gradient falls as $6/R^2$ along an opening path, and it can act as the second scale only if the coefficient in Eq.~\eqref{eq:a0} vanishes. These are negative statements in the scope of that dictionary. They point to the other chambers, without the Romans mass, and to other completions.

One comparison with an eleven-dimensional geometry survives as an observation. Witten's space with Killing spinors, $M^{111}$, is the hypercharge quotient of the $\bar5$ of $SU(5)$, since the weights $2,-3$ of its K\"ahler quotient are $6Y$ on $d^c$ and $L$; the cone is Calabi-Yau because $\mathrm{tr}\,Y=0$, and $11=4+(2\cdot5-3)$ (\file{02}, \cite{Witten1981,Fabbri1999}). The $Z_6$ of the Standard Model is then the kernel of $SU(3)\times SU(2)$ on it, so no bosonic Kaluza-Klein mode is a triplet or a doublet. The weights are in \cite{Fabbri1999} and the sources it builds on; a targeted search found no statement of the hypercharge reading, and that search was not exhaustive.

\section{Where this leaves the question}

Three things seem settled, in their scopes. The finite Dirac operator, grading and real structure are boundary data of a parent, exactly. The Higgs field is the interval component of a connection, seen in the window. Noncommutative finite algebras arise from commutative parents as limits of compression maps, and a local operator can produce the window.

Three things stay open, in order of consequence. The first is the vector sector on a seven-dimensional parent: nothing yet identifies the tunnelling scale with $M_W$, and the coupling $g_4^{-2}=\mathrm{Vol}/g^2$ falls with the size. The second is the selection of the family, by a term in the action with a local dependence on the order parameter, or by a stabilized shape. The third is the part of the finite data that remains input: on the interval the chirality comes from the choice of walls, and on the disc bundle the $\mathbf 4$ is pseudoreal, so the chiral spectrum of the Standard Model is not obtained. Equivariant cobordism invariance of the index \cite{Witten1983} gives total index zero on the boundary, so a chiral spectrum needs another source: walls with data beyond the extension of a bundle, or fluxes on an even base.

\section*{Status}

The statements of the rows marked refereed were read by GPT-6.1 Terra or GPT-6 Astra, as the ledger records. Eq.~\eqref{eq:a0}, Proposition~\ref{prop:so30}, Eq.~\eqref{eq:shadow} and the corrections of the last section but one are new or newly assembled here, with proofs in \file{10} and \file{11}. They have had an independent reading by Claude Sonnet 5.5, which found an error in our first count of the gauge generators in Eq.~\eqref{eq:a0} and several amendments, all applied; they have had no referee of the kind the other rows had, and the corrections have not been read again. Numerical statements are Python, in the repository tools. Nothing here is submitted.

\section*{Provenance}

This note has no listed authors. It was written by Claude Opus 5.5. It draws on lemmas by Claude Opus 5.5 and GPT-6 Astra, on referee reports by GPT-6.1 Terra, on an audit by GPT-6 Astra, and on the eleven-dimensional line written by another Claude session.

\bibliography{towerless_intervals}

\end{document}
